\section{Additional Multi-Severity Analysis}
\label{sec:appendix-multiseverity}

\providecommand{\vTwoNumberMacros}{paper1_v2_high_severity_draft/tables/multiseverity_numbers}
\input{\vTwoNumberMacros}
\providecommand{\vTwoAssociationTable}{tables/table_multiseverity_associations}

This section gives the aggregation and sensitivity details behind
\Cref{sec:exp-multiseverity-extension}. The broad-severity experiment is
analyzed separately from the original narrow sweep. It keeps the original
thresholds and pair construction fixed and uses continuous retention rather
than defining a new performance band.

The same \vTwoVetoAtZeroEight{} TwoRoom checkpoints account for all SR failures at
$\sigma_{\mathrm{eval}}=\vTwoEvalZeroEight$ and
$\vTwoEvalZeroSixteen$: one has $\sigma_{\max}=\vTwoTrainOne$ and two have
$\sigma_{\max}=\vTwoTrainOneFive$. At
$\sigma_{\mathrm{eval}}=\vTwoEvalZeroThirtyTwo$, the failures comprise
\vTwoVetoThirtyTwoTwoRoom{} TwoRoom, \vTwoVetoThirtyTwoPushT{} PushT, and
\vTwoVetoThirtyTwoReacher{} Reacher checkpoints; at
$\vTwoEvalZeroSixtyFour$, they comprise
\vTwoVetoSixtyFourTwoRoom{} TwoRoom and \vTwoVetoSixtyFourPushT{} PushT
checkpoints. No Cube checkpoint passes IR while failing SR. At the higher
severities, the failures span TwoRoom, PushT, and Reacher, although some
checkpoint rows recur across severity slices and the conditional denominator
also changes with severity.

For each task and training run, the grid contains eight training-noise maxima
and four evaluation severities. Let $x_{ij}$ denote any diagnostic or behavioral
quantity at training level $i$ and evaluation severity $j$. We remove the two
additive main effects using
\[
\widetilde{x}_{ij}
=x_{ij}-\overline{x}_{i\cdot}-\overline{x}_{\cdot j}
+\overline{x}_{\cdot\cdot}.
\]
Spearman correlations are computed over the \vTwoFactorialCellCount{} centered
cells within a single task and training run. We then average the
\vTwoRunCount{} training-run correlations within each task and weight the
\vTwoTaskCount{} tasks equally. The task ranges in
\Cref{tab:multiseverity-associations} summarize heterogeneity; they are not
confidence intervals. Double-centering leaves
$(\vTwoTrainLevelCount-1)(\vTwoEvalLevelCount-1)=\vTwoInteractionDf$
interaction degrees of freedom and imposes row- and column-sum constraints, so
the \vTwoFactorialCellCount{} cells are not treated as independent observations.

\input{\vTwoAssociationTable}
\FloatBarrier

Using stress success divided by same-checkpoint clean success instead of their
percentage-point difference gives the same qualitative result. The equal-task
mean direct correlations are $\vTwoIRRatioDirect$ for lower relative IR and
$\vTwoSRRatioDirect$ for higher SR; after accounting for the exposure ratio
they are $\vTwoIRRatioAdjusted$ and $\vTwoSRRatioAdjusted$, respectively.
Every clean-success estimate is nonzero (minimum
$\vTwoMinimumCleanSuccess\%$), but ratios can still amplify changes for weak
checkpoints; we therefore keep the percentage-point difference as the primary
outcome and treat the ratio only as a sensitivity check. These are descriptive
factorial associations. The \vTwoRowCount{} rows are not independent
replicates, and the four evaluation-severity rows for a checkpoint share the
same clean-performance estimate; therefore we do not attach row-level
$p$-values or confidence intervals.
